
Large language models can generate code that passes a high proportion of benchmark tests, yet their success rate drops when given their own compiler errors and asked to repair the code. This asymmetry between code generation capability and self-repair capability suggests that the representation of error information may affect model performance.

Prior work has established that including compiler feedback in self-repair loops improves pass rates. Chen et al. (2025) report minimum +4.9 percentage point improvement across model scales. However, existing approaches typically feed raw compiler output directly to models without modifying its format. The assumption that any error message format is equally useful to the model has not been systematically tested.

Compiler errors are formatted for terminal display with terse, technical syntax. LLMs are trained primarily on natural language prose, documentation, and code with explanatory comments. This distributional mismatch between compiler output format and model training distribution motivates the hypothesis that reformatting errors toward natural language structure could improve model comprehension.

This work tests whether organizational structure---independent of information content---affects self-repair success. The key methodological contribution is a scrambled control condition that contains identical diagnostic information as structured format but with randomly permuted section order. If structured format outperforms scrambled format, the effect can be attributed to organizational structure rather than information availability.

The contributions of this work are:

\begin{enumerate}
\item \textbf{Evidence for structure effects}: Structured error formatting outperforms scrambled formatting with identical content (+14.4 percentage points, p < 0.001), indicating that organizational structure affects repair success independently of information content.
\end{enumerate}

\begin{enumerate}
\item \textbf{Information preservation validation}: A reconstruction test achieves 100\% accuracy across 500 samples, confirming that format transformation preserves all diagnostic information.
\end{enumerate}

\begin{enumerate}
\item \textbf{Fix specificity pattern}: Simulation results show an inverted-U relationship between hint specificity and repair success, with intermediate hints (Level 2) outperforming both no hints and exact fixes.
\end{enumerate}

\begin{enumerate}
\item \textbf{Scale interaction null finding}: While a directional pattern exists (smaller models show larger benefits), the Format $\times$ Model Scale interaction does not reach statistical significance.
\end{enumerate}
